\subsection{Structure of the Clifford algebra.}

In this section, we will assume that A is a field, that E is a vector space of finite dimension m over A, and that the quadratic form Q is non-degenerate (which, according to th. 1 of § 5, no. 1, requires that m be even if A has characteristic 2). Since E is free, the canonical map ρ is injective (no. 3, cor. 2 of th. 1). We shall henceforth identify E with its image in C(Q).

THEOREM 2. --- Suppose that the dimension of E is an even number \(m = 2r\) and that Q be neutral ( \(\S 4, \text{n}^{\circ} 2\) ). Then the algebra C(Q) is separable (Chap.~VIII, \(\S 7, \text{n}^{\circ} 5\) , Def. 1) and is isomorphic to the algebra of endomorphisms of a vector space of dimension \(2^{r}\) on A. Moreover, if \(m > 0\) , C \(^{+}\) (Q) is separable and is the direct sum of two ideals isomorphic to the algebra of endomorphisms of a vector space of dimension \(2^{r-1}\) over A.

Indeed, since Q is neutral, one can decompose E as a direct sum of two totally singular subspaces N and P of dimension r (§ 4, no. 2, cor. 1 of prop. 2). Since the restriction of Q to N is zero, the subalgebra S of C(Q) generated by N is identified with the exterior algebra of N (no. 3, lemma 4). For \(n \in \mathbb{N}\) , we shall denote by \(e_n'\) the map \(t \to nt\) from S into itself.

Let \((n_{1},\ldots,n_{r})\) be a basis of N; we shall denote by \((p_{1},\ldots,p_{r})\) the basis of P such that \(\Phi(n_{i},p_{j})=\delta_{ij}\) ( \(\S 4\) , no. 2, prop. 2). For \(p\in P\) , we shall denote by \(p'\) the linear form \(n\to\Phi(n,p)\) on N, and \(i_{p}\) the endomorphism of S obtained by passing to the quotient from the endomorphism \(i_{p'}\) of T(N) associated with \(p'\) as was stated in Lemma 1 of No. 2. By (5), we have

\[
e _ {n} ^ {\prime} \circ i _ {p} + i _ {p} \circ e _ {n} ^ {\prime} = \Phi (n, p) \quad (n \in \mathrm{N}, p \in \mathrm{P}).\tag{12}
\]

Let us set, for \(x = n + p \in \mathbf{E}\) (with \(n \in \mathbf{N}\) and \(p \in \mathbf{P}\) ), \(s(x) = e_n' + i_p\) . It is clear that \(s\) is a linear map from \(\mathbf{E}\) into \(\mathfrak{L}(\mathbf{S})\) . Since we have

\[
s (x) ^ {2} = \left(e _ {n} ^ {\prime} + i _ {p}\right) ^ {2} = Q (n) + \Phi (n, p) = Q (x)
\]

By virtue of (12) and Lemma 1 (no. 2), s extends to a homomorphism (which we shall still denote by s) from C(Q) into ℒ(S) (no. 1, prop. 1). We shall show that this homomorphism is surjective, which, since C(Q) and ℒ(S) are both of dimension 2²ʳ, will imply that s is an isomorphism and will prove our first assertion.

Indeed, let I denote the interval \([1, r]\) . For every subset H of I, we set \(H' = I - H\) and we denote by \(n_{H}\) (resp. \(p_{H}\) ) the product of \(n_{i}\) (resp. \(p_{i}\) ) for \(i \in H\) , arranged in increasing order of indices. Recall that the \(n_{H}\) form a basis of S (no. 3, th. 1). Finally, for any two subsets H, K of I, set \(x_{H,K} = n_{H}p_{I}n_{K'}\) . We shall show that the elements \(s(x_{H,K})\) of \(s(C(Q))\) generate L(S). Now, if \(j \notin H\) , one has \(s(p_{j})(n_{H}) = i_{p_{j}}(n_{H}) = 0\) by lemma 1, since the \(n_{i}\) for \(i \in H\) belong to the kernel of the linear form \(n \to \Phi(n, p_{j})\) on N; on the other hand we have

\[
s \left(p _ {j}\right) \left(n _ {j} n _ {\mathrm{H}}\right) = \left(i _ {p _ {j}} \circ e _ {n _ {j}} ^ {\prime}\right) \left(n _ {\mathrm{H}}\right) = \Phi \left(p _ {j}, n _ {j}\right) n _ {\mathrm{H}} - n _ {j}. s \left(p _ {j}\right) \left(n _ {\mathrm{H}}\right) = n _ {\mathrm{H}}
\]

(by (12)). Since \(s\) is a homomorphism, we deduce, for any two subsets H, K of I, that \(s(p_{\mathrm{K}})(n_{\mathrm{H}}) = 0\) if \(\mathrm{K} \not\subset \mathrm{H}\) , and that \(s(p_{\mathbf{K}}) (n_{\mathbf{H}}) = \pm n_{\mathbf{H} - \mathbf{K}}\) if \(\mathrm{K} \subset \mathrm{H}\) . Since, for \(\mathrm{M} \subset \mathrm{I}\) and \(\mathrm{L} \subset \mathrm{I}\) , we have by definition \(s(n_{\mathbf{M}})(n_{\mathbf{L}}) = n_{\mathbf{M}}n_{\mathbf{L}}\) , and that \(n_{\mathbf{M}}n_{\mathbf{L}}\) is zero if \(\mathrm{M} \cap \mathrm{L} \neq \emptyset\) and is equal to \(\pm n_{\mathbf{M} \cup \mathbf{L}}\) in the contrary case, we conclude from the preceding that, for arbitrary subsets \(\mathrm{H}\) , \(\mathrm{K}\) , \(\mathrm{L}\) of \(\mathrm{I}\) , \(s(x_{\mathbf{H},\mathbf{K}})(n_{\mathbf{L}}) = s(n_{\mathbf{H}})s(p_{\mathbf{I}})s(n_{\mathbf{K}^{\prime}})(n_{\mathbf{L}})\) is zero if \(\mathrm{K} \neq \mathrm{L}\) and is equal to \(\pm n_{\mathbf{H}}\) if \(\mathrm{K} = \mathrm{L}\) . This shows that the \(s(x_{\mathbf{H},\mathbf{K}})\) generate \(\mathfrak{L}(\mathrm{S})\) and completes the proof of the first assertion.

To prove the second assertion, set \(S^{+} = S \cap C^{+}\) and \(S^{-} = S \cap C^{-}\) ; it is clear that \(S^{+}\) (resp. \(S^{-}\) ) is the subspace of S generated by the \(n_{\mathrm{H}}\) such that H has an even (resp. odd) number of elements, that S is the direct sum of \(S^{+}\) and \(S^{-}\) , and that \(s(C^{+})\) leaves \(S^{+}\) and \(S^{-}\) stable. Consequently \(s\) maps \(C^{+}\) into a subalgebra of \(\mathfrak{L}(S)\) , isomorphic to the product \(\mathfrak{L}(S^{+}) \times \mathfrak{L}(S^{-})\) ; the restriction of \(s\) to \(C^{+}\) is an isomorphism of \(C^{+}\) onto this subalgebra, since \(s\) is injective and \(C^{+}\) and \(\mathfrak{L}(S^{+}) \times \mathfrak{L}(S^{-})\) both have dimension \(2^{2r-1}\) (no. 2, cor. 1 of th. 1). QED.

COROLLARY. --- If m is even, but Q of arbitrary index, the algebra C(Q) is a central simple algebra of dimension \(2^{m}\) . Moreover, if m \textgreater{} 0, the subalgebra C⁺(Q) is separable, and its center Z has dimension 2 over A. When Z is a field, Z is a separable quadratic extension of A and C⁺(Q) is simple; otherwise, Z is a direct product of two fields isomorphic to A, and C⁺(Q) is then a direct product of two simple subalgebras of dimensions \(2^{m-2}\) .

Indeed, let A' be the algebraic closure of A, and Q' the quadratic form on \(E' = A' \otimes_{A} E\) obtained from Q by extension of scalars. We have seen that \(C(Q')\) is isomorphic to \(A' \otimes_{A} C(Q)\) (prop. 2), and it is clear that \(C^{+}(Q')\) is isomorphic to \(A' \otimes_{A} C^{+}(Q)\) . Since Q' is neutral (§ 4, no. 2, cor. 2 of prop. 3), the corollary is an immediate consequence of th. 2 and the permanence theorems of chap.~VIII, § 7.

Remarks. --- 1) Since the algebra C(Q) is simple, it has only one class of irreducible representations; these are called the spinor representations; when one fixes one's attention on one of these representations, say \(\tau\) , the elements of the space in which \(\tau\) takes place are called spinors. If Q is neutral, the restriction of \(\tau\) to \(C^{+}(Q)\) is, like that of \(s\) , sum of two inequivalent absolutely irreducible representations; the elements of the subspaces carrying these two representations are called semi-spinors. In the general case, if \(C^{+}(Q)\) is not simple, the restriction of \(\tau\) to \(C^{+}(Q)\) must, since it is faithful, contain subrepresentations belonging to each of the two classes of irreducible representations of \(C^{+}(Q)\) , and hence is a sum of two inequivalent absolutely irreducible representations, since this is the case after extending scalars to the algebraic closure A' of A. On the other hand, if \(C^{+}(Q)\) is simple, it has only one class of irreducible representations, and the restriction of \(\tau\) to \(C^{+}(Q)\) is therefore irreducible, since it decomposes under scalar extension to A' into two non-equivalent representations.

\begin{enumerate}
\def\labelenumi{\arabic{enumi})}
\setcounter{enumi}{1}
\tightlist
\item
  Suppose A has characteristic \(\neq 2\) , and let \((x_{1},\ldots,x_{m})\) ( \(m=2r\) ) be an orthogonal basis of E. Put
\end{enumerate}

\[
z = 2 ^ {r} x _ {1} \dots x _ {m} \in \mathrm{C(Q)};
\]

since \(x_{i}x_{j} + x_{j}x_{i} = 0\) for \(i \neq j\) , one has \(zx_{j} = -x_{j}z\) , which implies that \(z\) belongs to the center \(Z\) of \(C^{+}(Q)\) without belonging to \(A\) . We have

\[
z ^ {2} = 2 ^ {2 r} (- 1) ^ {r} \mathrm{Q} (x _ {1}) \dots \mathrm{Q} (x _ {m}) = (- 1) ^ {r} \mathrm{D}
\]

denoting by D the discriminant of \(\Phi\) with respect to the basis \((x_{j})\) (cf.~exerc. 9).

THEOREM 3. --- Suppose that the dimension of E is an odd number \(m = 2r + 1\) (so that A has characteristic \(\neq 2\) ).

\begin{enumerate}
\def\labelenumi{\alph{enumi})}
\item
  The algebra C \(^{+}\) (Q) is central simple. If Q has maximum index r, C \(^{+}\) (Q) is isomorphic to the algebra of endomorphisms of a vector space of dimension 2 \(^{r}\) over A.
\item
  The algebra C(Q) is separable. Its center Z has dimension 2, and C(Q) is isomorphic to Z⊗ \(_{A}\) C \(^{+}\) (Q), hence is simple or a direct product of two simple subalgebras.
\end{enumerate}

Indeed, let \(x_0\) be a non-isotropic vector of E, and F the orthogonal complement of \(x_0\) ; let us denote by \(Q_1\) the quadratic form \(y \to -Q(x_0)Q(y)\) on F; it is clear that \(Q_1\) is nondegenerate. Since \(x_0y = -yx_0\) (for \(y \in F\) ), we have \((x_0y)^2 = -Q(x_0)Q(y) = Q_1(y)\) , and consequently the map \(y \to x_0y\) of F into \(C^+(Q)\) extends to a homomorphism \(h\) of \(C(Q_1)\) in \(C^+(Q)\) (no. 1, prop. 1). But \(C(Q_1)\) is simple (th. 2) and has the same dimension \(2^{2r}\) as \(C^{+}(Q)\) ; this implies, since \(h(1) = 1\) , that \(h\) is an isomorphism. Moreover, if \(Q\) has index \(r\) , one can choose \(x_0\) in such a way that \(Q_1\) also has index \(r\) ( \(\S 4\) , no. 2, prop. 3), which proves \(a\) ).

Now let \((x_{1},\ldots ,x_{2r})\) be an orthogonal basis of F; set \(z = x_0x_1\dots x_{2r}\) . One immediately verifies that \(z\) commutes with \(x_{j}\) for \(j = 0,\dots,2r\) , hence belongs to the center of C(Q). Let Z be the subspace of C(Q) generated by 1 and \(z\) ; it is a subalgebra of the center of C(Q) and a quadratic extension of A, since \(z\) is odd and \(z^2\) is equal to the scalar \((-1)^r\mathrm{Q}(x_0)\ldots \mathrm{Q}(x_{2r})\) . Consider the homomorphism \(\theta\) of \(\mathbf{Z}\otimes_{\mathbf{A}}\mathbf{C}^{+}(\mathbf{Q})\) into C(Q) defined by \(\theta (u\otimes v) = uv\) . Since \(z\in \mathbb{C}^{-}\) and is invertible, the map \(u\to zu\) is an isomorphism of the module \(\mathbf{C}^+\) onto \(\mathbb{C}^{-}\) , which implies that \(\theta (\mathrm{Z}\otimes \mathrm{C}^{+})\) contains \(\mathbf{C}^{+}\) and \(\mathbb{C}^{-}\) , hence coincides with C(Q). Since \(\mathbf{Z}\otimes \mathbf{C}^{+}\) and C(Q) have the same dimension \(2^{2r + 1},\theta\) is an isomorphism; this proves \(b)\) taking into account the results of Chap.~VIII, § 7.

Note. --- The discriminant D of \(\Phi\) with respect to the basis \((x_{j})_{(j=0,\ldots,2r)}\) is equal to \(2^{2r+1}Q(x_{0})\ldots Q(x_{2r})\) . Consequently, Z is generated by 1 and by the odd element \(z' = 2^{r+1}z\) such that \(z'^{2} = (-1)^{r}2D\) . The algebra C(Q) is therefore simple if and only if \(2(-1)^{r}D\) is not a square in A.
% semantic-fragments: legacy-input-seam
\relax{}
